% Yuri Y. Vieira Sugano — currículo (PT, orientado a web/dados, uma página).
% Versão em português do résumé EN. Header com links em texto (compila em qualquer lugar).
\documentclass[a4paper,10pt]{article}
\usepackage[utf8]{inputenc}
\usepackage[T1]{fontenc}
\usepackage[brazil]{babel}
\usepackage[top=0.6in, bottom=0.6in, left=0.9in, right=0.9in]{geometry}
\usepackage{enumitem}
\usepackage{titlesec}
\usepackage{xcolor}
\usepackage[colorlinks=true, linkcolor=lightblue, urlcolor=lightblue]{hyperref}
\usepackage{graphicx}
\graphicspath{{icons/}}   % coloque github.png, linkedin.png, site.png em yuri/resume/icons/

\definecolor{darktext}{RGB}{0, 51, 102}
\definecolor{lightblue}{RGB}{66, 139, 202}

\titleformat{\section}{\large\bfseries\color{darktext}}{}{0em}{}[\titlerule]
\titlespacing{\section}{0pt}{6pt}{3pt}
\setlist[itemize]{leftmargin=1.2em, topsep=2pt, itemsep=1pt, parsep=0pt}
\setlength{\parindent}{0pt}

\begin{document}

\begin{center}
  {\LARGE \color{darktext} \textbf{Yuri Y. Vieira Sugano}}\\[2pt]
  \textit{R. Jaracatiá 305, São Paulo, Brasil \;|\; (11) 99647-6219 \;|\;
  \href{mailto:yuri@artelonga.com.br}{yuri@artelonga.com.br}}\\[2pt]
  \href{https://github.com/yurisugano}{\includegraphics[height=0.45cm]{github.png}} \quad
  \href{https://www.linkedin.com/in/yuri-sugano-a55304153/}{\includegraphics[height=0.45cm]{linkedin.png}} \quad
  \href{https://artelonga.com.br/}{\includegraphics[height=0.45cm]{site.png}}
\end{center}

\vspace{2pt}
Engenheiro de software e dados, e neurocientista. Constrói e publica serviços web full-stack sozinho
(Rust/Axum, SvelteKit, Node), com sólida experiência em dados/ML (Python, Spark, PLN). Todos os sistemas em produção.

\section*{Experiência}
\textbf{Fundador} — \href{https://artelonga.com.br/}{ArteLonga} \hfill 2024--atual\\
Serviços web para pequenos negócios --- 6 deployments em produção (Rust/Node na Fly.io).\\[3pt]
\textbf{Quant} — PointSet Technologies \hfill 2024--2026\\
Sistemas quantitativos e inteligência para mercados financeiros.\\[3pt]
\textbf{Senior Data Scientist} — Propz \hfill 2023--2024\\
Liderança de projetos e arquitetura técnica para analytics de varejo/negócios em sistemas de big data
(Apache Spark, Delta Lake); alinhando prioridades de negócio à entrega.\\[3pt]
\textbf{Técnico de Pesquisa} — Universidade de Chicago, Neurobiologia \hfill 2018--2023\\
Aquisição, processamento e interpretação de dados; gestão de equipe e projetos; ensino e comunicação científica.\\[3pt]
\textit{Antes:} Monitor de Estatística, U.\ de Chicago (2019--2020); Gerente de Laboratório, U.\ de Chicago
(2015--2017); Assistente de Pesquisa, Instituto Butantan — bolsista FAPESP (2012--2013).

\section*{Formação}
\textbf{Bacharelado em Neurociência} (concentração em Química), Universidade de Chicago, 2018--2022 ---
GPA 3,7/4,0. Foco: neurociência do comportamento social.

\section*{Habilidades}
\textbf{Linguagens:} Python, R, JavaScript/TypeScript, Rust, Scala, SQL.\\
\textbf{Dados \& ML:} Apache Spark, Delta Lake, scikit-learn, TensorFlow, spaCy, UMAP, DBSCAN, OpenCV.\\
\textbf{Web \& Infra:} SvelteKit, Axum, Node; Docker, Fly.io, Git, Jenkins, Azure; AWS (S3, EC2, EMR, QuickSight), GCP.\\
\textbf{Ferramentas:} Markdown, LaTeX, Vim, Obsidian, Quartz.

\section*{Portfólio Técnico}
\href{https://co.artelonga.com.br}{\textbf{co}} --- plataforma de conteúdo colaborativo, sync em tempo real
\hfill \textit{Rust/Axum, SvelteKit, LiteFS}\\
\href{https://rfq.artelonga.com.br}{\textbf{rfq}} --- motor de precificação RFQ em tempo real (Avellaneda-Stoikov)
\hfill \textit{Rust/Axum}\\
\href{https://yggdrasil.artelonga.com.br}{\textbf{yggdrasil}} --- host de jogos Godot 4 / WASM
\hfill \textit{Rust/Axum, WASM}\\
\href{https://github.com/yurisugano/SensorySpeech}{\textbf{SensorySpeech}} --- PLN para descrições
sensoriais \hfill \textit{Python, scikit-learn, spaCy, UMAP, DBSCAN}\\
\href{https://github.com/yurisugano/ProprioSuite}{\textbf{ProprioSuite}} --- pacote para testes
proprioceptivos \hfill \textit{Python, TensorFlow, SQL, OpenCV}\\
\href{https://github.com/yurisugano/Bystander-Effect-Rat-Analysis}{\textbf{Bystander Effect in Rats}}
--- análise estatística de comportamento social \hfill \textit{R, mineração de dados}

\section*{Publicações \& Patente}
\begin{itemize}
  \item \href{https://doi.org/10.1101/2022.07.01.498150}{Vieira Sugano, Y.\,Y., et al.\ (2022). Helping can be motivated by non-affective cues in rats. \textit{bioRxiv}}.
  \item \href{https://doi.org/10.1126/sciadv.abb4205}{Havlik, J.\,L., Vieira Sugano, Y.\,Y., et al.\ (2020). The bystander effect in rats. \textit{Science Advances}}.
  \item Madelaire, C.\,B., Vieira Sugano, Y.\,Y., et al.\ (2020). Challenges of dehydration in invasive toads. \textit{Behavioral Ecology and Sociobiology}.
  \item \textbf{Patente:} \href{https://patents.google.com/patent/US11589798B2/en}{Theseometer for Measuring Proprioception Performance} — US 11.589.798 B2, Universidade de Chicago.
\end{itemize}

\vspace{2pt}\textit{\small Bibliografia completa \& referências: \href{https://neuro.artelonga.com.br/2026-05-29/}{neuro.artelonga.com.br/2026-05-29}.}

\end{document}
